\section{Conclusiones}

Un launch file de simulación multi-robot combina cuatro patrones que trabajan en conjunto:

\textbf{OpaqueFunction} permite que el número de nodos a lanzar sea variable y dependa de argumentos de la terminal. Sin él, el launch file solo puede declarar nodos de forma estática.

\textbf{IncludeLaunchDescription} permite reutilizar los launch files oficiales de Gazebo y RSP sin duplicar su contenido. El launch padre construye la secuencia de alto nivel; los hijos se encargan de sus detalles.

\textbf{TimerAction} resuelve las dependencias temporales entre componentes. Gazebo necesita tiempo antes de recibir spawns; el RSP necesita tiempo antes de que el spawn lea su tópico; el nodo de control necesita que el robot exista antes de intentar leer odometría.

\textbf{Namespaces} permiten que múltiples instancias del mismo nodo coexistan sin interferencia. Es el único mecanismo para separar los canales de comunicación de distintos robots sin modificar el código del nodo.

Los Scripts~\ref{lst:consenso_part1}--\ref{lst:consenso_part3} muestran el launch file completo del experimento dividido en sus tres bloques principales:

{\renewcommand{\lstlistingname}{Script}
\begin{figure*}[t]
\centering
\begin{lstlisting}[style=python,
    caption={consenso\_prom.launch.py -- Parte 1: imports y poses iniciales.},
    label={lst:consenso_part1}]
import os
import math

from launch import LaunchDescription
from launch.actions import (
    DeclareLaunchArgument, IncludeLaunchDescription,
    OpaqueFunction, TimerAction
)
from launch.launch_description_sources import PythonLaunchDescriptionSource
from launch.substitutions import LaunchConfiguration
from launch_ros.actions import Node
from ament_index_python.packages import get_package_share_directory

# Posiciones (x, y, yaw) de los 8 robots en el mundo de Gazebo.
# El layout esta disenado para que todos los pares esten a menos
# de 5.0 m entre si, garantizando conectividad con ese radio.
POSES_8 = [
    (0.0, 0.0,   math.pi / 4),      # robot0 -- esquina BL
    (4.0, 1.0,  -math.pi / 2),      # robot1 -- borde inferior
    (8.0, 0.0,  -math.pi / 4),      # robot2 -- esquina BR
    (1.0, 4.0,   math.pi),          # robot3 -- borde izquierdo
    (8.0, 4.0,   0.0),              # robot4 -- borde derecho
    (0.0, 8.0,   3*math.pi / 4),    # robot5 -- esquina TL
    (4.0, 9.0,   math.pi / 2),      # robot6 -- borde superior
    (8.0, 8.0,  -3*math.pi / 4),    # robot7 -- esquina TR
]
\end{lstlisting}
\end{figure*}
}

{\renewcommand{\lstlistingname}{Script}
\begin{figure*}[t]
\centering
\begin{lstlisting}[style=python,
    caption={consenso\_prom.launch.py -- Parte 2: construccion dinamica de acciones.},
    label={lst:consenso_part2}]
def launch_setup(context, *args, **kwargs):

    # Leer y convertir los argumentos ya resueltos por el contexto
    n              = int(LaunchConfiguration('n_robots').perform(context))
    R              = float(LaunchConfiguration('neighbor_radius').perform(context))
    t_final        = float(LaunchConfiguration('t_final').perform(context))
    spawn_obstacle = LaunchConfiguration('spawn_obstacle').perform(context).lower() == 'true'
    obs_x          = float(LaunchConfiguration('obstacle_x').perform(context))
    obs_y          = float(LaunchConfiguration('obstacle_y').perform(context))

    poses         = POSES_8[:n]
    simulador_pkg = get_package_share_directory('simulador')
    obstacle_sdf  = os.path.join(simulador_pkg, 'models', 'obstacle_cylinder.sdf')
    articubot_pkg = get_package_share_directory('articubot_one')
    gazebo_launch = os.path.join(get_package_share_directory('gazebo_ros'),
                                 'launch', 'gazebo.launch.py')
    rsp_launch    = os.path.join(articubot_pkg, 'launch', 'rsp.launch.py')

    actions = []

    # 1. Gazebo arranca de inmediato para tener el mayor tiempo de
    #    inicializacion antes de que lleguen los primeros spawns.
    actions.append(
        IncludeLaunchDescription(PythonLaunchDescriptionSource(gazebo_launch))
    )

    # 2. Obstaculo opcional: se spawnea a los 3 s si el argumento lo pide.
    if spawn_obstacle:
        actions.append(TimerAction(period=3.0, actions=[
            Node(package='gazebo_ros', executable='spawn_entity.py',
                 arguments=['-file', obstacle_sdf, '-entity', 'obstacle_cylinder',
                             '-x', str(obs_x), '-y', str(obs_y), '-z', '0.5'],
                 output='screen')
        ]))

    # 3. Por cada robot: RSP -> spawn -> control, con 1.5 s entre robots
    #    para no saturar Gazebo durante el arranque masivo.
    for i, (x0, y0, yaw) in enumerate(poses):
        ns    = f'robot{i}'
        delay = 3.0 + i * 1.5

        rsp_node = IncludeLaunchDescription(
            PythonLaunchDescriptionSource(rsp_launch),
            launch_arguments={'namespace': ns, 'use_sim_time': 'false'}.items()
        )
        spawn_node = Node(
            package='gazebo_ros', executable='spawn_entity.py',
            arguments=['-topic', f'/{ns}/robot_description', '-entity', ns,
                       '-x', str(round(x0,4)), '-y', str(round(y0,4)),
                       '-z', '0.1', '-Y', str(round(yaw,4))],
            output='screen'
        )
\end{lstlisting}
\end{figure*}
}

{\renewcommand{\lstlistingname}{Script}
\begin{figure*}[t]
\centering
\begin{lstlisting}[style=python,
    caption={consenso\_prom.launch.py -- Parte 3: timers, nodos tardios y punto de entrada.},
    label={lst:consenso_part3}]
        control_node = Node(
            package='control_nodes', executable='consenso_node',
            name='consenso', namespace=ns,
            parameters=[{'consensus_type': 'prom_err',
                         'k1': 0.8, 'k2': 1.0, 'vmax': 1.0, 'wmax': 1.0}],
            output='screen'
        )
        actions.append(TimerAction(period=delay,       actions=[rsp_node]))
        actions.append(TimerAction(period=delay + 0.5, actions=[spawn_node]))
        actions.append(TimerAction(period=delay + 1.5, actions=[control_node]))

    # 4. Nodos de analisis: esperan a que el ultimo robot este operativo.
    late_delay = 3.0 + n * 1.5 + 2.0
    actions.append(TimerAction(period=late_delay, actions=[
        Node(package='simulador', executable='aire', name='aire',
             parameters=[{'neighbor_radius': R}], output='screen')
    ]))
    actions.append(TimerAction(period=late_delay, actions=[
        Node(package='simulador', executable='states_plotter_v2',
             parameters=[{'save_dir': 'figuras/consenso_prom',
                          't_final': t_final, 'n_robots_expected': n}],
             output='screen')
    ]))
    actions.append(TimerAction(period=late_delay, actions=[
        Node(package='simulador', executable='pva_plotter',
             parameters=[{'save_dir': 'figuras/consenso_prom', 't_final': t_final}],
             output='screen')
    ]))
    return actions


# Punto de entrada: declara los argumentos configurables
# y delega la construccion a launch_setup via OpaqueFunction.
def generate_launch_description():
    return LaunchDescription([
        DeclareLaunchArgument('n_robots',        default_value='8'),
        DeclareLaunchArgument('neighbor_radius', default_value='5.0'),
        DeclareLaunchArgument('t_final',         default_value='60.0'),
        DeclareLaunchArgument('spawn_obstacle',  default_value='false'),
        DeclareLaunchArgument('obstacle_x',      default_value='4.0'),
        DeclareLaunchArgument('obstacle_y',      default_value='4.5'),
        OpaqueFunction(function=launch_setup)
    ])
\end{lstlisting}
\end{figure*}
}

